% Глава 2. Проектирование клиентской части платформы.
% Источники: knowledge-base/05-design-system.md,
%            knowledge-base/06-architecture-overview.md,
%            knowledge-base/personal-contribution.md.

\clearpage
\section{Проектирование клиентской части платформы}

\subsection{Архитектура экосистемы <<ВКурсе>>}

\subsubsection{Состав клиентских приложений}
% TODO: контент
Клиентская часть платформы <<ВКурсе>> построена как набор автономных приложений, объединённых общим серверным контуром и единым продуктовым подходом. В её состав входят: мобильное приложение для iOS, мобильное приложение для Android, веб-приложение, включающее лендинг и публичный просмотр расписания, и закрытый веб-сервис администрирования для сотрудников учебного заведения. Все четыре компонента работают с единым серверным API и при этом сохраняют независимость в части стека технологий и UX-решений, обеспечивая нативность опыта на каждой платформе.

\subsubsection{Взаимодействие с бэкенд-сервисом}
% TODO: контент
Взаимодействие клиентских приложений с серверной частью построено по схеме <<запрос --- преобразование --- доставка>>. На текущем этапе платформа работает в двух режимах: парсинговом (бэкенд извлекает расписание с публичных источников шести подключённых омских вузов и преобразует его в единый формат) и собственном (бэкенд хранит и валидирует данные, поступившие через веб-сервис администрирования). С каждым запросом клиента передаётся уникальный идентификатор устройства, формируемый при первом запуске приложения; это позволяет вести подсчёт уникальных и новых пользователей без сбора персональных данных. Подробности контракта раскрываются в смежной выпускной квалификационной работе, посвящённой серверной части платформы; в настоящей работе серверная часть рассматривается только в объёме, необходимом для описания клиентских компонентов.

\subsection{Проектирование пользовательского опыта}

\subsubsection{Принцип нативности и единого подхода}
% TODO: контент
В основе продуктового подхода к проектированию пользовательского опыта лежит принцип максимальной нативности: каждое клиентское приложение следует рекомендациям соответствующей платформы и использует её родные компоненты, паттерны навигации и визуальный язык. Этот подход осознанно противопоставлен распространённой практике построения единой кросс-платформенной дизайн-системы поверх различных гайдлайнов. Аргументы в пользу нативного подхода: единообразие с системными приложениями, узнаваемость взаимодействий для пользователя, отсутствие необходимости компенсировать различия платформ собственными решениями.

\subsubsection{Информационная архитектура}
% TODO: контент
Информационная архитектура клиентов опирается на общую модель данных, в которой центральными сущностями являются учебное заведение, группы, преподаватели, аудитории и собственно расписание. На уровне пользовательского интерфейса каждое клиентское приложение предоставляет согласованный набор сценариев: выбор основной сущности (группа или преподаватель), просмотр актуального и предстоящего расписания, работа с избранным, заметки к парам и персонализация.

\subsubsection{Дизайн-система и Figma-файлы}
% TODO: контент
Для каждой из трёх платформ поддерживается отдельный Figma-файл; все файлы хранятся в единой Figma-организации проекта. Такая декомпозиция отражает принцип нативности: попытка ведения общего файла означала бы выбор <<общего знаменателя>> между Apple Human Interface Guidelines, Material Design~3 Expressive и shadcn/ui, что не давало бы оптимального результата ни на одной из платформ.

\subsection{Проектирование под особенности платформ}

\subsubsection{iOS: Human Interface Guidelines, адаптация под iOS 18\textendash{} и iOS 26+}
% TODO: контент
Проектирование iOS-приложения опирается на Apple Human Interface Guidelines. Особое внимание уделено двойной адаптации: на iOS 18 и ниже используются классические системные компоненты, на iOS 26 и выше задействуется визуальный язык Liquid Glass во всех уместных элементах интерфейса. Эта дифференциация выполнена средствами SwiftUI и условной маршрутизации компонентов по версии операционной системы.

\subsubsection{Android: Material Design 3 Expressive}
% TODO: контент
Android-приложение спроектировано на основе Material Design 3 Expressive. В реализации используется динамическая темизация с поддержкой пользовательских акцентных цветов, что согласуется с экспрессивным характером актуальной версии Material и обеспечивает естественный вид приложения в системе.

\subsubsection{Web: shadcn/ui и адаптивная вёрстка}
% TODO: контент
Веб-часть проекта построена на компонентах shadcn/ui поверх Tailwind CSS. Этот выбор обеспечивает современный визуальный язык, согласованность лендинга, публичного веб-просмотра расписания и закрытого сервиса администрирования, а также позволяет владеть исходным кодом UI-компонентов и адаптировать их под продуктовые требования. Вёрстка адаптивна и одинаково корректно отображается на десктопе и на мобильных устройствах.

\subsection{Проектирование ключевых сценариев}

\subsubsection{Сценарий студента: просмотр и персонализация}
% TODO: контент
Базовый сценарий студента состоит из выбора собственной группы (одноразово при первом запуске), последующего ежедневного обращения к актуальному расписанию и точечной персонализации (избранные группы и преподаватели, заметки к парам, выбор темы и акцентного цвета). Сценарий спроектирован так, чтобы повседневное использование требовало минимального числа действий: главный экран сразу показывает текущую или ближайшую пару.

\subsubsection{Сценарий методиста: внесение изменений и распространение расписания}
% TODO: контент
Сценарий методиста реализуется в веб-сервисе администрирования. Сотрудник учебного отдела работает с централизованной моделью данных, вносит изменение в конкретный слот расписания, в реальном времени получает предупреждения о коллизиях (занятая аудитория, занятый преподаватель, конфликт по группе) и подтверждает изменение. После сохранения обновлённое расписание автоматически распространяется во все клиентские интерфейсы; повторной публикации на сторонних ресурсах не требуется.
